\section{结论}\label{sec:conclusion}
本文针对未知宽带强弱重叠源，将波形消元与位置搜索结合，形成多候选联合AP和波形消元块NOMP。前者保留可能的起点，再联合调整两源位置；后者在位置试探时重估波形，减少固定系数带来的重复小步。

2400条观测中，五候选MCJ-AP达到99\%的双源20\,m成功率，相对单候选恢复96条，平均耗时0.718\,s。相同五候选配对中，联合精修保持定位表现，将平均耗时从1.084\,s降至0.744\,s，减少31.36\%，梯度计算减少42.30\%。构型分组与候选案例表明，多候选改善困难区域的恢复，联合精修降低分支处理成本。

VP-BNOMP在全部2400条上达到94.92\%成功率，平均耗时0.586\,s。同批共同240条中，消元更新保持固定系数版本的成功数，将平均耗时从2.942\,s降至0.591\,s，与曲率分析相符。两种流程分别强调候选覆盖与局部更新效率；后续可扩展至未知源数和更一般的传播条件。
